% Written by Nghi (PL) — Project Lead
\section{Threats to Validity}
\label{sec:threats}

\subsection{Internal Validity}

\textbf{Researcher degrees of freedom.}
The research questions, hypotheses, threshold (75\%), statistical tests, and
multiplicity-correction scheme (Holm) were registered in the proposal prior to
data collection.
The analysis script was finalised before the full run began.
Three protocol deviations were recorded as formal amendments in \texttt{notes.md}
before any metric was inspected, preventing post-hoc hypothesis adjustment.

\textbf{Model non-determinism.}
We pin the exact model snapshot (\texttt{gpt-4o-mini-2024-07-18}) and set
\texttt{temperature}\,=\,0 to minimise output variance.
We do not claim complete determinism because the provider does not guarantee it
for parallel decoding; however, a single response per unit was collected and
stored verbatim, making the study a single-point observation rather than a
probabilistic claim.

\textbf{Infrastructure retries.}
Network-level retries were permitted only for transport failures (HTTP 5xx before
content delivery).
After the API returned any content, no re-call was issued.
This preserves the first-attempt intent of the study.

\subsection{Construct Validity}

\textbf{Branch coverage as quality proxy.}
Branch coverage does not capture semantic correctness of test oracles.
We mitigate this by using mutation score (PIT) as a co-primary metric for Java,
which directly assesses the ability of tests to detect injected logical faults.

\textbf{Python mutation score unavailable.}
\texttt{pytest-mutagen 1.3} requires explicit mutant decorators rather than
automatic injection; consequently, zero mutants were generated for 20 of 22
executable Python units, and the remaining two aborted because the base test suite
failed on the unmutated code.
Python mutation score is therefore excluded from confirmatory analysis.
This limits the comparability of quality evidence between the two languages.

\textbf{Non-executable tests.}
One Java unit (\texttt{parseCsvLine}, CC\,=\,9) and three Python units produced
non-executable tests.
All four are retained in the denominator; our primary analysis uses complete-case
(executable units only), and a worst-case sensitivity analysis assigns 0\% to
non-executable units.
Both analyses agree in their rejection of $H_0$ for RQ1.

\subsection{External Validity}

\textbf{Small and filtered sample.}
Our sample consists of 50 standalone functions with $5 \le \text{CC} \le 15$,
drawn from a single frozen repository.
Results should not be generalised beyond this complexity range or to
class-level, file-level, or framework-dependent code.
Effect sizes and 95\% confidence intervals are reported throughout to aid
readers in assessing practical significance.

\textbf{Single model snapshot.}
All results are conditioned on \texttt{gpt-4o-mini-2024-07-18}.
Performance of other models or future snapshots of the same model may differ.
RQ3 (CC vs.\ quality correlation) is marked exploratory to caution against
over-generalising the complexity-coverage relationship.

\textbf{Randoop fairness.}
GPT-4o-mini and Randoop operate under identical conditions: same Java source
revision, same JDK 17, same PIT toolchain, and a fixed Randoop time budget
of 60 seconds per unit with a pinned random seed (\texttt{20260621}).
The time budget was chosen to reflect a realistic CI-integration scenario;
longer budgets may narrow the observed gap.

\textbf{Tool and language heterogeneity.}
Java coverage is measured by JaCoCo and Python coverage by coverage.py;
the two tools may differ in their branch-counting granularity for equivalent
control-flow constructs.
Confirmatory RQ1 tests are conducted separately for each language; pooled results
are used only for descriptive purposes.
